\section*{الفصل \ref{ch:g10:chemical-species} --- الأنواع الكيميائية الطبيعية والاصطناعية}

\begin{solution}{exo:g10:chemical-species:1}
طبيعية: الفانيلين من القرن، والكافيين من حبوب البن. اصطناعية: الأسبرين،
وفانيلين المصنع.
\end{solution}

\begin{solution}{exo:g10:chemical-species:2}
لأنهما \omterm{def:g6:pure-substances-and-mixtures:species}{النوع الكيميائي} نفسه، وللنوع الخصائص نفسها أيًّا كان
مصدره.
\end{solution}

\begin{solution}{exo:g10:chemical-species:3}
$R_f = 2.4/6.0 = 0.40$.
\end{solution}

\begin{solution}{exo:g10:chemical-species:4}
يبرّد البخار ويعيده سائلًا. ودخول الماء من الطرف السفلي يجعله يملأ الغلاف كله
(ويلاقي أسخن البخار في الأخير، فيكون تبريده
أفضل).
\end{solution}

\begin{solution}{exo:g10:chemical-species:5}
أن يذيب النوع أفضل بكثير من الماء؛ وأن يكون \omterm{def:g6:solutions-and-solubility:miscible}{غير قابل للامتزاج} بالماء؛ وأن
يكون قليل الخطر قدر الإمكان وسهل الإزالة.
\end{solution}

\begin{solution}{exo:g10:chemical-species:6}
الإيثانول \omterm{def:g6:solutions-and-solubility:miscible}{قابل للامتزاج} بالماء: فلا تتكوّن طبقتان، ولا يمكن فصل شيء. أما
الهكسان الحلقي \omterm{def:g6:solutions-and-solubility:miscible}{فغير قابل للامتزاج} بالماء ويذيب اللينالول جيدًا.
\end{solution}

\begin{solution}{exo:g10:chemical-species:7}
الماء في الأعلى: لأن \qty{1}{mL} من ثنائي كلورو الميثان يزن \qty{1.33}{g}،
أكثر من \qty{1}{mL} من الماء، فيغوص.
\end{solution}

\begin{solution}{exo:g10:chemical-species:8}
ينصهر عند درجة أدنى من اللازم (الأسبرين: \qty{135}{\celsius}) وعلى مدى من
الدرجات: فالعيّنة غير نقية، أو ليست أسبرين.
\end{solution}

\begin{solution}{exo:g10:chemical-species:9}
اللينالول وإيثانوات اللينالين (فبقعتاه في ارتفاعي المرجعين). وإيثانوات
اللينالين صعد أعلى.
\end{solution}

\begin{solution}{exo:g10:chemical-species:10}
\omterm{def:g10:chemical-species:distillation}{التقطير المائي}؛ ثم إضافة الهكسان الحلقي إلى المقطَّر والرج؛ ثم فصل الطبقتين؛
ثم \omterm{def:g3:separating-mixtures:evaporate}{تبخير} الهكسان الحلقي.
\end{solution}

\begin{solution}{exo:g10:chemical-species:11}
رصاص القلم لا \omterm{def:g2:mixing-and-dissolving:dissolve}{يذوب} في \omterm{def:g10:chemical-species:tlc}{المُزيح}؛ وفوق \omterm{def:g10:chemical-species:tlc}{المُزيح} تُحمل البقع صعودًا في الصفيحة بدلًا
من أن تذوب في الحوض.
\end{solution}

\begin{solution}{exo:g10:chemical-species:12}
يغلي الهكسان الحلقي عند \qty{81}{\celsius}، واللينالول عند \qty{198}{\celsius}:
فالتسخين الهادئ يطرد الهكسان الحلقي بينما يبقى اللينالول.
\end{solution}

\begin{solution}{exo:g10:chemical-species:13}
العيّنة على الأرجح ذلك النوع (فلها $R_f$ \omterm{def:g10:chemical-species:natural}{النوع الطبيعي} نفسه على الصفيحة
نفسها) وعلى الأرجح نقية (بقعة واحدة، ودرجة انصهار حادة).
\end{solution}

\begin{solution}{exo:g10:chemical-species:14}
$0.40 \times 5.5 = \qty{2.2}{cm}$. والآخر عند $0.75 \times 5.5 =
\qty{4.1}{cm}$ (إلى أقرب مليمتر)، أي فوق A بنحو \qty{1.9}{cm}.
\end{solution}

\begin{solution}{exo:g10:chemical-species:15}
يستطيع \qty{200}{mL} أن يذيب $1.6 \times 0.200 = \qty{0.32}{g}$: فمع
\qty{0.20}{g} ليس مشبعًا؛ ويستطيع أن يذيب \qty{0.12}{g} أخرى.
و\qty{0.20}{g} الذائبة كانت ستضيع مع الماء؛ أما الهكسان الحلقي، الذي يذيب
اللينالول أفضل بكثير، فيستعيد معظمها.
\end{solution}

\begin{solution}{pb:g10:chemical-species:1}
\textbf{1.} بخاره يحمل الأنواع العطرية من الأزهار إلى
المكثّف.

\textbf{2.} يبرّد البخار فيعيده سائلًا؛ ويدخل الماء البارد من الطرف
السفلي.

\textbf{3.} \omterm{def:g6:solutions-and-solubility:miscible}{غير قابل للامتزاج}؛ وأقل كثافة من الماء (فهو يطفو).

\textbf{4.} الطبقة رقيقة جدًا، وجزء من الزيت ذائب في الماء أو منتشر فيه:
\omterm{def:g10:chemical-species:extraction}{والاستخلاص} يستعيده.

\textbf{5.} يذيب اللينالول جيدًا، \omterm{def:g6:solutions-and-solubility:miscible}{وغير قابل للامتزاج} بالماء، ويغلي عند درجة
منخفضة (\qty{81}{\celsius})، فإزالته سهلة.

\textbf{6.} لا: فالإيثانول يمتزج بالماء ولا تتكوّن طبقتان.

\textbf{7.} في الأعلى (\qty{0.78}{g} لكل mL، أقل من الماء). ومع ثنائي
كلورو الميثان (\qty{1.33}{g} لكل mL) تكون طبقة \omterm{def:g6:solutions-and-solubility:solute}{المذيب} في
الأسفل.

\textbf{8.} لا لهب بالقرب (قابل للاشتعال)؛ والعمل تحت ساحبة الغازات
بالقفازات، وجمع النفايات بدلًا من صبها في المغسلة (سام للأحياء
المائية).

\textbf{9.} يغلي الهكسان الحلقي عند \qty{81}{\celsius}، واللينالول عند
\qty{198}{\celsius}: فالتسخين الهادئ لا يزيل إلا \omterm{def:g6:solutions-and-solubility:solute}{المذيب}.

\textbf{10.} $R_f(\text{L}) = 2.4/6.0 = 0.40$؛ $R_f(\text{A}) = 4.5/6.0
= 0.75$.

\textbf{11.} اللينالول وإيثانوات اللينالين.

\textbf{12.} لا: فهو يحوي نوعين على الأقل؛ إنه \omterm{def:g6:pure-substances-and-mixtures:mixture}{خليط}.

\textbf{13.} لكي تُنشر البقع كلها في الشروط نفسها: فعندئذ فقط يمكن مقارنة
الارتفاعات.

\textbf{14.} قيمة $R_f$ له تساوي قيمة اللينالول في الشروط نفسها: فهو على
الأرجح لينالول، مطابق \omterm{def:g10:chemical-species:natural}{للنوع الطبيعي}.

\textbf{15.} $1.0 \times 0.90 = \qty{0.90}{g}$.

\textbf{16.} $0.90/100 = \qty{0.90}{\percent}$.

\textbf{17.} $100/0.009 \approx 11\,000$\,g، أي نحو \qty{11}{kg} من
الأزهار.

\textbf{18.} فهو مستخلص من نبات، دون \omterm{def:g5:irreversible-changes:chemical-change}{تغير كيميائي}. لكن لينالوله هو النوع
نفسه الذي في اللينالول الاصطناعي: الصيغة نفسها، والخصائص نفسها.

\textbf{19.} \qty{0.90}{g} من الزيت العطري من \qty{100}{g} من الأزهار.
\end{solution}
